%==========================================================%
% preambule.tex — compléments au template TEK-UP            %
%                                                           %
% Le template charge déjà : xcolor[table], array, longtable, %
% multirow, enumitem, float, graphicx, caption, minitoc,     %
% biblatex, babel[french].                                   %
% Ce fichier ajoute ce qui manque aux chapitres 3 à 6.       %
%==========================================================%

%=== Mathématiques ========================================%
% Requis par le chapitre 4 (formules de stock, lissage
% exponentiel) et le chapitre 5 (WAPE, MASE) :
% environnements align, split, equation ; \sqrt, \Phi, \sum.
\usepackage{amsmath}
\usepackage{amssymb}

%=== Type de colonne des tableaux =========================%
% Colonne de largeur fixe alignée à gauche, utilisée par
% tous les longtable des chapitres.
% (Le template définit déjà L, R et C sans argument.)
\newcolumntype{P}[1]{>{\raggedright\arraybackslash}p{#1}}

%=== Bibliographie ========================================%
% Le template pointe vers biblio.bib ; le présent rapport
% utilise references.bib (44 entrées, style IEEE).
\addbibresource{references.bib}

%=== Réglages de flottants ================================%
% Les chapitres 3 à 6 contiennent beaucoup de figures pleine
% largeur ; on assouplit les seuils de placement.
\renewcommand{\topfraction}{0.9}
\renewcommand{\bottomfraction}{0.8}
\renewcommand{\textfraction}{0.07}
\renewcommand{\floatpagefraction}{0.75}
\setcounter{topnumber}{3}
\setcounter{bottomnumber}{2}
\setcounter{totalnumber}{5}

%=== Césure ===============================================%
% Le template désactive la césure (package hyphenat[none]).
% Certains tableaux contiennent des termes longs : on
% autorise un léger débordement plutôt qu'une ligne cassée.
\tolerance=2000
\emergencystretch=2em
